\section{Introduction}

Forensic analytics can prioritize investigation. They do not prove fraud,
certify outcomes, or replace risk-limiting audits. That boundary comes first
because residual methods are easy to overread: a large residual can be a useful
triage signal, but it is not a finding about intent, legality, or correctness.

This paper covers two W-series analytic families. Turnout residuals ask whether
a reporting unit's turnout is unusual after comparison with its own history,
nearby or similar units, and public covariates such as population composition or
registration patterns. Candidate-share residuals ask whether a contest share is
unusual relative to comparable units, previous contests, or cast vote record
features where those records are available. In both cases, a residual means
observed minus model-predicted: the reported value for a unit less the value a
baseline model would have expected from the comparison population.

The separation from the V-series is structural. V-series work verifies canvass
and audit claims under explicit evidence contracts. W-series residual reports
produce an investigation queue. They may help an analyst decide which precinct,
batch, scanner, or contest deserves review, but they cannot certify the election
and must not be stored or read as certification evidence. The W.01 overview
fixes that claim boundary and the common forensic report contract
\citep{wseries-overview-2026}; this paper instantiates the contract for turnout
and vote-share residuals.
